\begin{theorem}[Dirichlet 扭曲 Fourier 级数的收敛半径恰由谱半径决定，故平方根解析窗必然破裂]\label{thm:xi-output-potential-dirichlet-fourier-radius-barrier}
\leanverified{paper\_xi\_output\_potential\_dirichlet\_fourier\_radius\_barrier}
沿用命题 \ref{prop:real-input-40-output-dirichlet-nonrh} 的记号。固定整数 \(m\ge 2\) 与 \(1\le j\le m-1\)，记
\[
M_j:=M(\omega_m^j),
\qquad
\rho_{m,j}:=\rho(M_j),
\qquad
a_n(\omega_m^j):=\Tr(M_j^n).
\]
定义普通生成函数
\[
A_{m,j}(z):=\sum_{n\ge 1}a_n(\omega_m^j)\,z^n.
\]
则 \(A_{m,j}(z)\) 为有理函数，且其收敛半径恰为
\[
\boxed{\;R_{m,j}=\frac{1}{\rho_{m,j}}.\;}
\]
特别地，若某个模数 \(m\) 满足
\[
\theta_m:=\max_{1\le j\le m-1}\frac{\log\rho_{m,j}}{\log\lambda}>\frac12,
\qquad
\lambda:=\rho(M(1))=\varphi^2,
\]
则存在 \(j\) 使
\[
R_{m,j}<\lambda^{-1/2}.
\]
等价地，归一化级数
\[
A_{m,j}(\lambda^{-1/2}z)
\]
在单位圆内部已出现极点。因而不存在统一的平方根型解析窗能够同时覆盖全部非平凡 Dirichlet 扭曲模态。
\end{theorem}
\begin{proof}
在代数闭包上把 \(M_j\) 的特征值按代数重数记为
\[
\eta_1,\dots,\eta_{40}.
\]
则对每个 \(n\ge 1\) 都有
\[
a_n(\omega_m^j)=\Tr(M_j^n)=\sum_{\ell=1}^{40}\eta_\ell^n.
\]
故在 \(|z|<1/\rho_{m,j}\) 内可逐项求和，得到
\[
\begin{aligned}
A_{m,j}(z)
&=
\sum_{n\ge 1}\sum_{\ell=1}^{40}\eta_\ell^n z^n\\
&=
\sum_{\ell=1}^{40}\sum_{n\ge 1}(\eta_\ell z)^n\\
&=
\sum_{\ell=1}^{40}\frac{\eta_\ell z}{1-\eta_\ell z}.
\end{aligned}
\]
因此 \(A_{m,j}\) 为有理函数。其全部非零极点正位于
\[
z=\eta_\ell^{-1}
\qquad(\eta_\ell\neq 0).
\]
若某个非零特征值 \(\eta\) 的代数重数为 \(r\ge 1\)，则对应分式项合并为
\[
\frac{r\,\eta z}{1-\eta z},
\]
其留数非零，故该极点不会被消去。于是 \(A_{m,j}\) 的最近极点到原点的距离恰为
\[
\min_{\eta_\ell\neq 0}\frac{1}{|\eta_\ell|}
=
\frac{1}{\max_\ell|\eta_\ell|}
=
\frac{1}{\rho_{m,j}},
\]
这就证明了
\[
R_{m,j}=\frac{1}{\rho_{m,j}}.
\]

现若 \(\theta_m>1/2\)，则由定义存在某个 \(j\) 使
\[
\rho_{m,j}=\lambda^{\theta_{m,j}}
>
\lambda^{1/2}.
\]
于是
\[
R_{m,j}=\rho_{m,j}^{-1}<\lambda^{-1/2}.
\]
把变量按 \(\lambda^{-1/2}\) 归一化后，级数 \(A_{m,j}(\lambda^{-1/2}z)\) 的收敛半径变为
\[
\lambda^{1/2}R_{m,j}=\frac{\lambda^{1/2}}{\rho_{m,j}}<1.
\]
由于该函数为有理函数，收敛半径严格小于 \(1\) 等价于单位圆内部已有极点。故任何试图以 \(\lambda^{n/2}\) 作为统一尺度的解析窗，在该 Dirichlet 模态上都必然失败。证毕。
\end{proof}

\endinput
